\section{Introduction}
\label{sec:introduction}

Software maintenance frequently requires engineers to extract an existing
capability from a large system: a parser must become a small library, a
configuration resolver must be reused in another service, or a plugin registry
must be separated from the framework that originally hosted it. This operation
is common in modernization, dependency reduction, service decomposition, and
reuse of legacy code. It is also deceptively difficult. The target behavior may
be distributed across helper functions, stateful registries, resources,
configuration defaults, error classes, and third-party dependencies. Copying
the most obvious implementation region may produce a package that works on a
happy path while silently changing edge behavior or retaining hidden
dependencies on the original repository.

Repository-level code agents appear well suited to this task. They can search
large workspaces, inspect call sites and tests, edit multiple files, execute
tools, and package new artifacts. However, current evaluation paradigms do not
isolate whether an agent can perform behavior-preserving feature extraction.
Issue-resolution benchmarks ask an agent to repair a repository until tests
pass~\cite{jimenez2024swebench}. Greenfield code-generation benchmarks ask for
a new implementation from a specification~\cite{chen2021codex}.
Repository-understanding and localization studies test whether the relevant
code can be found. None of these settings simultaneously requires an agent to
recover an existing feature from an entangled repository, preserve a declared
behavioral contract, remove dependence on the original package, and avoid
solving the task through broad repository copying.

Feature extraction therefore combines at least three obligations:
\begin{enumerate}
  \item \textbf{Locate:} identify implementation evidence relevant to the
        requested capability;
  \item \textbf{Close:} recover the APIs, helpers, state, resources,
        dependencies, and edge semantics needed for behavioral completeness;
  \item \textbf{Isolate:} package the recovered feature so that it executes
        independently of the source repository.
\end{enumerate}
We use this decomposition as an explanatory model rather than a causal
factorization. A final failure does not by itself reveal which obligation
failed; semantic labels require trajectory and artifact evidence.

We introduce \flb, a benchmark that evaluates these obligations under a
controlled Full-Repository / No-Hint protocol. Each task provides a complete
repository pinned to an immutable source revision and a complete public
contract describing the required package interface and observable behavior.
The agent may inspect the entire repository, but it is not given source-location
hints, benchmark tests, Hidden tests, or a reference solution. It must produce
a standalone \texttt{featurelifted} package. A submission passes only when it
builds, satisfies both primary and deeper contract tests, and remains functional
after access to the original repository is removed.

\flb separates two quality dimensions. \emph{Functional Pass} measures whether
the submission satisfies build, Public, Hidden, and isolation gates.
\emph{Reference-Relative Extraction Size} (\rres) measures the normalized
footprint of a functionally passing package relative to an evaluator-side
feasible reference. Correctness and compactness are not combined into one
score: a broad vendoring solution may preserve behavior but fail to demonstrate
a compact extraction, while a small package may simply be incomplete.

The main Python-200$'$ suite combines a frozen 150-task baseline with a
separately calibrated Hard-50 split, for 200 tasks from 176 repositories. The
Hard-50 split increases coverage of plugin registries, lifecycle and session
behavior, multi-source configuration, validation boundaries, deep parsing, and
direct copy-trap tasks. An earlier External-50 expansion is retained only as an
easy, copy-heavy side split: strong-agent pass rates of 90--94\% and
pass-conditioned footprints close to whole-repository copying showed that it
was unsuitable for the paper's main difficulty claim.

Our empirical study asks four groups of questions:
\begin{itemize}
  \item \textbf{Capability:} How often do current code agents produce
        independent, behavior-complete feature packages?
  \item \textbf{Quality:} When agents pass, how compact are their extractions?
  \item \textbf{Failure mechanisms:} At which gate do agents fail, and which
        contract obligations remain open?
  \item \textbf{Information and resources:} How do executable feedback,
        process scaffolds, and trajectory cost change outcomes?
\end{itemize}

Our analysis points to a recurring gap between plausible repository work and
verifiable feature closure. Strong agents frequently produce buildable packages
and pass primary behavior checks, yet fail on required exports, boundary
conditions, exception semantics, stateful behavior, resource closure, or
preservation details. We refer to this recurring symptom as
\emph{contract-closure failure}. It is not inferred from a Hidden failure alone:
attribution requires evidence that the relevant requirement was public and that
the trajectory or artifact omitted or changed it. Controlled information
experiments further suggest that executable Public feedback and Hidden
behavioral completeness are distinct layers.

This paper makes three main contributions:
\begin{itemize}
  \item \textbf{Benchmark asset.} We define repository-level,
        behavior-preserving feature extraction and construct a suite of 200
        tasks from 176 pinned Python repositories, with public contracts,
        source registries, deterministic evaluators, evaluator-side feasible
        references, isolation checks, and task-level provenance. A v2
        contract-completeness label on freeze~v2 is 200/0/0 after repairing
        32 confirmed surface or entry-point defects recorded on the
        predecessor freeze (then 168/32/0), without changing suite membership.
  \item \textbf{Evaluation protocol.} We introduce a Full-Repository / No-Hint
        protocol that fixes repository, contract, runtime, budget, and evaluator
        while varying model backends. We report deterministic Functional Pass
        separately from pass-conditioned compactness and process cost.
  \item \textbf{Diagnostic analysis.} We analyze failure stages, semantic
        contract-closure symptoms, task factors, information boundaries, and
        trajectory cost. We also report negative results from legal
        self-testing, repair, checkpoint, context, and budget-control
        interventions rather than presenting them as successful methods.
\end{itemize}

\draftnote{After Gate A and Gate C, add one sentence with the eligible
Python-200-prime capability range and the dominant failure-stage result.}

